\scanpage{18}
\mspage{7}
% Folio d'archive : 17.
pour $\forall A$, ce sera O.K.

Voici un cas où il en sera ainsi : Supposons que $f$ admette
un adjoint à droite (disons) $g$ satisfaisant aux conditions
\[
\begin{array}{ll}
\text{(i)} & f^{-1}(W')=W \\[3pt]
\text{(ii)} & fg(X')\longrightarrow X'\ \in W'
\qquad \forall\unsure{X|X'}\in\operatorname{Ob}M'.
\end{array}
\]
Alors
\diagram{\begin{tikzcd}[column sep=huge]
\uline{\operatorname{Hom}}(\unsure{A|M},M)
\arrow[r,shift left=1.1ex,"{f^{\inserted{M}}_{\deleted{A}}}"] &
\uline{\operatorname{Hom}}(A,M')
\arrow[l,shift left=1.1ex,"{g^{\inserted{M}}_{\deleted{A}}}"]
\end{tikzcd}}
sont également adjoints l'un de l'autre, et \deleted{\unsure{les}}
satisfont $=^{\mathrm{s}}$ conditions pour $W_A,W'_A$. Donc
\diagram{\begin{tikzcd}[column sep=huge]
\mathcal H_W(A)\arrow[r,shift left=1.1ex,"{f_W^A}"] &
\mathcal H_{W'}(A)\arrow[l,shift left=1.1ex,"{g_W^A}"]
\end{tikzcd}}
% Inscription d'une utilisation antérieure de la feuille, à 180 degrés,
% barrée de plusieurs traits horizontaux, chevauchée par le diagramme.
\deleted{Monsieur Grothendieck}

Les théories $\mathcal H$ qui peuvent se décrire par un $(M,W)$ sont dites
\uline{spéciales}. Cela implique des conditions draconiennes sur $(M,W)$
